\section{Related Work}\label{sec:related}

This section places our work among six lines of research. For each line it states what others have done, how
our work relates to it and what we reuse. It closes with the gaps that remain open and with the data each
line was evaluated on, and Table~\ref{tab:related} summarizes both.

\subsection{Program slicing and srcSlice}
Weiser introduced program slicing~\cite{weiser1984program} and showed that programmers use slices when they
debug~\cite{weiser1982programmers}. Dependence graphs reduced slicing to graph
reachability~\cite{ottenstein1984pdg, ferrante1987pdg, horwitz1990interprocedural}, Podgurski and Clarke gave
dependence a semantic definition~\cite{podgurski1990formal}, and dynamic slicing restricts a slice to one
execution~\cite{korel1988dynamic, agrawal1990dynamic}. Surveys catalog the many kinds of slices that
followed~\cite{tip1994a, silva2012vocabulary, binkley2004a}. Probabilistic slicing preserves the distribution of the
criterion variables~\cite{hur2014slicing, amtoft2020theory, olmedo2025slicing}, and Gallagher names quantum
programs as a direction for semantically defined slicing~\cite{gallagher2025program}.

srcSlice computes slices from srcML, the XML form of source code that keeps every character and position of
the original~\cite{collard2011lightweight, collard2013srcml}. Instead of building a dependence graph, it records
for every variable a \emph{slice profile}, the lines where the variable is defined and used, the variables it flows
into and the calls it is passed to, and it computes dependences from these profiles as it slices. It slices the
Linux kernel in about 20 minutes with accuracy comparable to CodeSurfer~\cite{alomari2014srcslice}. General
frameworks such as WALA provide dependence-graph slicers for Python~\cite{wala}, but on a Qiskit program every
gate is a method call on one circuit object, so neither they nor srcSlice can tell which qubits an operation
depends on.

\emph{Relation and reuse.} We keep the classical notions, a criterion, a backward and a forward slice, and
dependence as the possibility of change, and define them for the reduced state of a set of qubits, a criterion
stronger than a distribution in one basis, which misses phase. The \tool{} approach reuses srcML's
representation and srcSlice's slice profiles for the classical part of a program and adds a quantum layer,
which Section~\ref{sec:approach} describes. We fix two gaps in srcSlice's handling of Python, subscript uses and
module globals, as a patch to the open-source tool.

\subsection{srcML for quantum programs}
Behler et al.\ specify a srcML markup for OpenQASM~3, so that quantum syntax shares one representation with the
classical languages srcML already supports~\cite{behler2024supporting}. On it they build an execution-order
analysis that unrolls loops, expands gate and function calls and labels operations inside \texttt{if}
statements with their condition, and their QStatic tool uses the resulting per-qubit order to detect
iteration and recursion patterns and to suggest refactorings~\cite{behler2025static}. Their analyzer does not
yet handle \texttt{while} loops or qubits passed to functions as classical values~\cite{behler2024supporting},
and its order table marks the control of a CNOT apart from its target. Behler et al.\ name, as future work, the
adaptation of srcML's slicer to find every line that affects a given qubit~\cite{behler2025static}.

\emph{Relation and reuse.} Our work takes that step. We reuse the idea that a static execution order of
qubit operations, kept on srcML, is the basis of quantum analyses, and we add what a slicer needs on top of it,
a dependence rule with a guarantee. An order table serves pattern detection well, but read directly as a
dependence graph it treats the control as read-only, which Section~\ref{sec:core} shows to be unsound. Their
analyzer reads OpenQASM~3 and ours reads Qiskit and OpenQASM~2, so the two cover different languages.

\subsection{Dependence graphs for quantum programs}
Yu and Zhao's quantum program dependence graph (QPDG)~\cite{yu2025qpdg} extends the classical
PDG~\cite{ferrante1987pdg} with entanglement, no-cloning, temporal and measurement dependences and proposes
slicing on it. It is defined in prose with examples and has not been implemented or evaluated. Its own example
shows why validation matters. In the graph Yu and Zhao draw for their first listing, the slice of a measured
qubit omits the two gates that fix its outcome, so the slice returns 1 with probability
\numYzFigSliceProb{} where the program returns 1 with certainty. Kaul et al.\ build a code property graph for
Qiskit and OpenQASM programs and query it for bugs and smells, among them superfluous operations, gates that
do not affect a measured qubit~\cite{kaul2023uniform}. Their graph treats the control of a CNOT as unchanged,
so the query can report gates whose effect reaches a measured qubit back through a control. Kaul et al.\
demonstrate the graph on example programs and note the lack of large datasets.
Roadmaps of quantum software engineering describe dependence graphs as an aid to understanding
programs~\cite{murillo2025roadmap}, and a roadmap of testing and debugging finds dependence and data-flow
analysis of circuits almost unexplored, naming only Kaul et al.~\cite{ramalho2025road}.

\emph{Relation and reuse.} We take the QPDG's question, which operations a qubit depends on, and answer it
against a semantic definition. We reimplement the QPDG as a documented baseline, logging every reading of its
prose before any run, and Section~\ref{sec:core} shows that its entanglement edges add no dependence the
write-all rule misses.

\subsection{Entanglement analysis}
Static analyses approximate which qubits may be entangled, with type systems~\cite{perdrix2007types,
prost2009reasoning, yuan2022twist}, abstract interpretation~\cite{perdrix2008entanglement, yu2021abstract}, the
stabilizer formalism~\cite{honda2015stabilizer} and data-flow rules~\cite{xia2025entanglement}. ScaffCC assumes
that qubits that interact may be entangled~\cite{javadiabhari2015scaffcc}, the rule our QPDG baseline uses.

\emph{Relation.} These analyses answer a different question, whether two qubits may be correlated, and they are
not slicers. Section~\ref{sec:core} shows that entanglement neither implies nor is implied by dependence, so
their results cannot decide what a slice keeps.

\subsection{Light cones and compilers}
Compilers build gate-dependence graphs to schedule and map circuits~\cite{murali2019triq, amy2020staq}, and
QAOA analyses bound the gates that can reach an observable by its light cone~\cite{farhi2014qaoa}. Qiskit's
\texttt{LightCone} pass keeps the gates that can reach an observable and drops those that commute with
it~\cite{qiskit2025}, and Heim et al.\ note that quantum dependences should account for
commutation~\cite{heim2020quantum}. Quantum separation logics reason locally about the qubits a program
touches~\cite{zhou2021qsl, le2022qsl}. Compilers also move a single-qubit gate within the idle time of its qubit
without changing the computation~\cite{smith2022timestitch}, and call the time windows of a circuit its
slices, a sense different from ours.

\emph{Relation and reuse.} The backward light cone is the relation of our write-all rule, and we do not claim
the relation as new. What these works do not provide is a definition of what a slice must preserve and a proof
that the light cone is exactly what a slice must keep, which Section~\ref{sec:core} gives. We reuse Qiskit's
\texttt{LightCone} pass as the commutation-pruning baseline.

\subsection{Debugging tools and bug detectors}
Metwalli and Van Meter adapt slicing to circuits by splitting a circuit at barriers the developer places and
removing the qubits a part does not use, and their tool maps gates to source lines through tracebacks at run
time~\cite{metwalli2022debugging, metwalli2024testing}. Their phase tests detect a wrong phase but need
hand-designed test vectors to locate the faulty qubit. Qiskit's slicing utilities partition a circuit by
barriers, depth or gate type~\cite{qiskitaddonutils}. A roadmap of testing and debugging lists slicing as a way
around no-cloning and notes that it grows hard to manage as circuits grow~\cite{ramalho2025road}, and testing
faces the related limit that phase faults escape computational-basis measurement~\cite{oldfield2024faster}.
Bug detectors find known patterns statically. LintQ runs CodeQL queries over Qiskit programs and orders two
operations when they share a qubit~\cite{paltenghi2024lintq}. QChecker matches patterns in the abstract syntax
tree~\cite{zhao2023qchecker}, and Q-PAC also detects known patterns~\cite{nayak2023qpac}. Bug
patterns~\cite{zhao2021bugpatterns}, Bugs4Q~\cite{zhao2023bugs4q}, mutation tools~\cite{mendiluze2021muskit,
fortunato2022qmutpy, garcia2022testing} and mutation-based fault localization~\cite{ishimoto2025mbfl} supply
faults and rankings to this field.

\emph{Relation and reuse.} Slices complement these tools. Manual slicing has no criterion, so the developer still
narrows the fault by hand, and a slice by the write-all rule generalizes the removal of unused qubits, since it
drops every qubit and operation that cannot reach the criterion. The \tool{} approach maps operations to source
lines statically. LintQ's ordering relation is the relation of the write-all rule, which confirms the rule in
practice, while LintQ uses it for patterns and makes no claim about slices. Detectors report known patterns,
and of the 42 Bugs4Q bugs on which QChecker was evaluated, 22 are wrong outputs that it cannot
detect~\cite{zhao2023qchecker}. A slice shows where such a fault can be. We reuse Bugs4Q, LintQ's corpus and its
public warnings, published mutation operators and mutation-based fault localization.

\subsection{Open gaps and the data used}
Three gaps remain open, and each can be shown from the work above. First, no dependence rule proposed for
quantum slicing has been checked against what a quantum program computes, and the one example worked out in
detail, the QPDG's, gives a wrong slice. Second, the analyses that represent quantum programs, from code
property graphs to execution-order tables, read the control of a gate as unchanged, and no work has measured
what that reading does to a slice. Third, the slicing support developers have is manual and dynamic, and no
static slicer shows a qubit's slice on Qiskit source with its loop iterations.

\begin{table*}[t]
  \centering
  \caption{Related approaches by what they analyze, whether they slice from a criterion, what they guarantee,
  whether they report source lines, and the data they were evaluated on.}
  \label{tab:related}
  \footnotesize\setlength{\tabcolsep}{3pt}
  \begin{tabular}{@{}>{\raggedright\arraybackslash}p{2.6cm} >{\raggedright\arraybackslash}p{2.4cm}
                  >{\raggedright\arraybackslash}p{2.6cm} >{\raggedright\arraybackslash}p{2.6cm} l
                  >{\raggedright\arraybackslash}p{4.4cm}@{}}
    \toprule
    Approach & Analyzes & Slicing criterion & Stated guarantee & Lines & Evaluated on \\
    \midrule
    srcSlice~\cite{alomari2014srcslice} & classical source, srcML & variable, automatic & none & yes & the Linux kernel, against CodeSurfer \\
    Behler et al., QStatic~\cite{behler2024supporting, behler2025static} & OpenQASM~3, srcML & none & none & yes & the 16 OpenQASM~3 example files \\ % TODO verify: QStatic's own evaluation set
    QPDG~\cite{yu2025qpdg} & quantum programs & variable or qubit & none & -- & examples, not implemented \\
    Kaul et al.~\cite{kaul2023uniform} & Qiskit, OpenQASM & none, queries & none & yes & example programs \\
    \texttt{LightCone}~\cite{qiskit2025} & circuits & observable & keeps the observable's value & no & -- \\
    Metwalli, Van Meter~\cite{metwalli2022debugging, metwalli2024testing} & Qiskit, at run time & none, manual breakpoints & none & yes, dynamic & small example circuits \\ % TODO verify: count of circuits and the survey
    LintQ~\cite{paltenghi2024lintq} & Qiskit, CodeQL & none, bug patterns & none & yes & 7,568 GitHub programs, \numLqTp{} confirmed bugs \\
    QChecker~\cite{zhao2023qchecker} & Qiskit, syntax tree & none, bug patterns & none & yes & the 42 bugs of Bugs4Q \\
    \midrule
    This paper, \tool{} & Qiskit, OpenQASM~2, srcML and srcSlice & a set of qubits, automatic & sound, dependence-complete & yes, static & \numNCircuits{} benchmark circuits, \numFamCircuits{} generated circuits, Bugs4Q, the Qiskit textbook, LintQ's \numScaleSOneFiles{} programs and warnings, QASMBench, Qiskit HumanEval \\
    \bottomrule
  \end{tabular}
\end{table*}

The data differ as much as the methods. The dependence graphs were shown on examples, the srcML work on the
OpenQASM~3 examples, and the detectors on Bugs4Q and on LintQ's GitHub corpus, while no slicing approach was
evaluated on real programs. We use the detectors' data where they fit, Bugs4Q for real faults and LintQ's corpus
and confirmed warnings for real programs, and add what a slicing study needs that they lack, circuits with a
simulated ground truth. These are \numNCircuits{} circuits from MQT Bench~\cite{quetschlich2023mqt}, QASMBench~\cite{li2023qasmbench} and
VeriQBench~\cite{chen2022veriqbench}, \numFamCircuits{} generated circuits with known structure, and seeded faults. For the scale of real programs we
add QASMBench's large circuits and, for current Qiskit, the expert-written programs of Qiskit
HumanEval~\cite{vishwakarma2024qhe}.
